% papex-template.tex -- Papex XeLaTeX 主文档
%
% 用途：配合 papex.cls 与 papex-build.py 使用。
%       构建器从 papex.json 读取元数据，生成若干 _papex_*.tex 片段与主文档一起编译。
%
% 作者只需：
%   1. 在下方「用户配置区」按需调整（通常无需改动）；
%   2. 撰写 sections/*.tex 章节源文件；
%   3. 运行 `python papex-build.py archive.tar.gz` 或 `make`。
%
% 注意：不要在此手动写 \title / \author —— 它们由 _papex_meta.tex 自动注入。
%       如需手动覆盖，可在 % ============================================================
% 本文件由 papex-build.py 自动生成，请勿手工编辑。
% 修改请在 papex.json 中进行，随后重新构建。
% ============================================================


\title{基于图神经网络的中文文献引用预测}
\papexSubtitle{一种融合语义与结构的端到端框架}
\papexLicense{CC-BY-4.0}
\papexVenue{第十二届中文计算与语言处理会议（CCLP 2026）}
\papexDoi{10.1234/papex.2026.00012}
\papexVersionNote{v1：初稿}
\papexKeywords{图神经网络；学术引用预测；中文文献；预训练语言模型；异构图注意力}
\papexRunningTitle{GNN 中文引用预测}
\author[1]{张明\thanks{通讯作者。Email: zhangming@example.edu.cn}\thanks{ORCID: 0000-0002-1825-0097}}
\author[1]{李华\thanks{同等贡献。}}
\author[2]{Wang Lei\thanks{同等贡献。}}
\affil[1]{山东大学 计算机科学与技术学院，济南}
\affil[2]{Institute of Computing Technology, CAS, Beijing}
\papexPdfAuthor{张明; 李华; Wang Lei}
 之后重新定义。

\documentclass[11pt]{papex}

% ===================== 用户配置区（可选） =====================
% 覆盖字体集（本地有系统中文字体时可改为 windows / mac / ubuntu）：
% \documentclass[11pt,fontset=fandol]{papex}
%
% 双栏排版：取消下一行注释
% \documentclass[11pt,twocolumn]{papex}
%
% 覆盖页面尺寸（需放在 \begin{document} 之前）：
% \geometry{margin=2.5cm}
%
% 额外宏包（按需）：
% \usepackage{algorithm,algorithmic}
% \usepackage{tikz}
% ============================================================

% ---- 自动生成区：由 papex-build.py 从 papex.json 注入 ----
% 这些文件在每次构建时重新生成；请勿手工编辑。
% ============================================================
% 本文件由 papex-build.py 自动生成，请勿手工编辑。
% 修改请在 papex.json 中进行，随后重新构建。
% ============================================================


\title{基于图神经网络的中文文献引用预测}
\papexSubtitle{一种融合语义与结构的端到端框架}
\papexLicense{CC-BY-4.0}
\papexVenue{第十二届中文计算与语言处理会议（CCLP 2026）}
\papexDoi{10.1234/papex.2026.00012}
\papexVersionNote{v1：初稿}
\papexKeywords{图神经网络；学术引用预测；中文文献；预训练语言模型；异构图注意力}
\papexRunningTitle{GNN 中文引用预测}
\author[1]{张明\thanks{通讯作者。Email: zhangming@example.edu.cn}\thanks{ORCID: 0000-0002-1825-0097}}
\author[1]{李华\thanks{同等贡献。}}
\author[2]{Wang Lei\thanks{同等贡献。}}
\affil[1]{山东大学 计算机科学与技术学院，济南}
\affil[2]{Institute of Computing Technology, CAS, Beijing}
\papexPdfAuthor{张明; 李华; Wang Lei}
        % \title, \author, \affil, \papexPaperId, \papexKeywords, ...
\addbibresource{references.bib}

\begin{document}

\maketitle

\begin{abstract}
% ============================================================
% 本文件由 papex-build.py 自动生成，请勿手工编辑。
% 修改请在 papex.json 中进行，随后重新构建。
% ============================================================

学术文献之间的引用关系构成了一张庞大的知识网络。准确预测一篇新论文可能引用的文献，对于文献推荐、科研趋势分析与学术检索具有重要意义。本文提出一种面向中文文献的图神经网络引用预测框架 Papex-GNN：首先利用预训练语言模型对中文标题与摘要进行语义编码，随后将其作为节点特征接入引文图，通过异构图注意力网络同时建模“主题相似”与“作者协作”两类边。在自建的中文论文数据集上，本文方法相较基于 TF-IDF 与协同过滤的基线，在 Recall@10 指标上提升了 12.4 个百分点。消融实验进一步表明，语义编码与结构建模两个模块均对最终效果有显著贡献。
    % 摘要正文（来自 papex.json 的 paper.abstract）
\end{abstract}

% ---- 正文：各章节由生成器按 papex.json 的 sections 顺序 \input ----
% ============================================================
% 本文件由 papex-build.py 自动生成，请勿手工编辑。
% 修改请在 papex.json 中进行，随后重新构建。
% ============================================================


\section{引言}
随着科研产出的指数级增长，如何从浩如烟海的文献中快速定位相关研究，成为科研工作者的核心痛点之一。引文网络（citation network）作为连接学术成果的结构化图谱，蕴含了丰富的“知识流向”信息：一篇论文的参考文献，本质上刻画了它在学术谱系中的来路。若能利用已有引文数据，预测一篇新论文最可能引用哪些文献，便能显著提升文献推荐与学术检索的体验。

传统方法主要依赖\emph{内容相似度}（如 TF-IDF + 余弦相似度）或\emph{协同过滤}（基于共同引用模式）\cite{kipf2017}。然而，前者难以捕捉深层语义，后者则面临严重的冷启动问题。近年来，图神经网络（GNN）在建模非欧结构数据上展现出强大能力\cite{wang2019}，为引文预测提供了新的建模工具；与此同时，以 Transformer 为代表的预训练语言模型\cite{vaswani2017,devlin2019} 能够为中文文本提供高质量的语义表示。

本文的主要贡献如下：

\begin{enumerate}
  \item 提出 Papex-GNN，一个面向中文文献的端到端引用预测框架，首次将预训练语义编码与异构图注意力结构建模统一在同一框架下；
  \item 设计了同时建模“主题相似边”与“作者协作边”的异构消息传递机制，显式区分两类结构信号；
  \item 在自建中文论文数据集上取得显著优于基线的效果，并提供了充分的消融分析。
\end{enumerate}

\section{相关工作}
\subsection{基于内容的引用预测}
早期工作将文献表示为词袋向量，通过相似度检索候选引用\cite{kipf2017}。这类方法实现简单，但受限于词汇鸿沟，难以识别表述不同但语义相近的工作。随着词嵌入与句向量的发展，语义匹配精度有所提升，但仍缺乏对全局引用结构的利用。

\subsection{基于图结构的引用预测}
将文献视为节点、引用关系视为边，引用预测可被形式化为链路预测问题。图卷积网络（GCN）\cite{kipf2017} 与图注意力网络（GAT）通过聚合邻居表征学习节点嵌入。然而，真实引文图是\emph{异构}的：除了引用边，还存在作者—论文、会议—论文等多种关系。异构图注意力网络（HAN）\cite{wang2019} 通过为不同元路径分配注意力权重，缓解了这一问题，但其注意力机制本身仍依赖输入节点的初始特征质量。

\subsection{预训练语言模型与中文处理}
以 BERT\cite{devlin2019} 与 ERNIE\cite{sun2020} 为代表的预训练模型，在中文自然语言理解任务上取得了突破性进展。本文采用这类模型的中文变体作为节点特征提取器，为图模型提供语义层面的初始化。

\section{方法}
\section{方法}
给定文献集合 $\mathcal{V}$ 与引用关系 $\mathcal{E}\subset \mathcal{V}\times\mathcal{V}$，我们的目标是学习一个打分函数 $s(u,v)$ 估计论文 $u$ 引用论文 $v$ 的概率。Papex-GNN 由语义编码与异构图聚合两个阶段组成。

\subsection{语义编码}
对每篇论文 $v$，将其标题与摘要拼接为文本 $t_v$，送入预训练中文编码器 $\mathrm{Enc}(\cdot)$，得到维度为 $d$ 的句向量：
\begin{equation}
  \mathbf{h}_v^{(0)} = \mathrm{Enc}(t_v) \in \mathbb{R}^{d}.
\end{equation}
该向量作为节点 $v$ 的初始特征，使图模型从第一步便具备语义判别力。

\subsection{异构图注意力聚合}
我们构建两类无向边：\textbf{主题相似边} $\mathcal{E}_{\mathrm{sim}}$（基于初始语义相似度动态构建）与\textbf{作者协作边} $\mathcal{E}_{\mathrm{co}}$（共享作者）。对每一类边 $r\in\{\mathrm{sim},\mathrm{co}\}$，采用注意力聚合：
\begin{equation}
  \alpha_{vu}^{r} = \frac{\exp\!\big(\mathrm{LeakyReLU}(\mathbf{a}_r^{\top}[\mathbf{W}_r\mathbf{h}_v \mathbin\Vert \mathbf{W}_r\mathbf{h}_u])\big)}
  {\sum_{u'\in\mathcal{N}_r(v)}\exp\!\big(\mathrm{LeakyReLU}(\mathbf{a}_r^{\top}[\mathbf{W}_r\mathbf{h}_v \mathbin\Vert \mathbf{W}_r\mathbf{h}_{u'}])\big)},
\end{equation}
其中 $\mathbf{W}_r$ 为类型相关变换矩阵，$\mathbf{a}_r$ 为注意力向量，$\mathcal{N}_r(v)$ 表示在关系 $r$ 下的邻居。最终节点表征为两类边表征的加权拼接：
\begin{equation}
  \mathbf{z}_v = \big\Vert_{r} \sigma\!\left(\sum_{u\in\mathcal{N}_r(v)} \alpha_{vu}^{r}\,\mathbf{W}_r\mathbf{h}_u\right).
\end{equation}

\subsection{训练目标}
采用成对排序损失（pairwise ranking loss），对每条真实引用 $(u,v^+)$ 采样若干负例 $v^-$：
\begin{equation}
  \mathcal{L} = \sum_{(u,v^+)} \sum_{v^-} \max\big(0,\, m - s(u,v^+) + s(u,v^-)\big),
\end{equation}
其中 $s(u,v)=\mathbf{z}_u^{\top}\mathbf{z}_v$，$m$ 为间隔超参数。

\section{实验}
\section{实验}
\subsection{数据集}
我们从公开中文学术平台收集了 $52\,148$ 篇论文及其引用关系，按 $8{:}1{:}1$ 划分训练、验证与测试集。统计信息如表~\ref{tab:stats} 所示。

\begin{table}[htbp]
  \centering
  \caption{数据集统计}
  \label{tab:stats}
  \begin{tabular}{lc}
    \toprule
    指标 & 数值 \\
    \midrule
    论文数 & $52\,148$ \\
    引用边数 & $318\,904$ \\
    作者数 & $76\,330$ \\
    平均引用数 & $6.12$ \\
    \bottomrule
  \end{tabular}
\end{table}

\subsection{基线方法}
我们对比了三类基线：(1) \textbf{TF-IDF}：基于词频—逆文档频率的余弦相似度；(2) \textbf{GCN}：仅用语义初始化特征的同构图卷积；(3) \textbf{HAN}：异构图注意力网络（随机初始化特征）。

\subsection{结果与分析}
主结果如表~\ref{tab:main} 所示。Papex-GNN 在 Recall@10 上达到 $0.487$，较最强的基线 HAN 提升 $12.4$ 个百分点。值得注意的是，GCN 在引入语义初始化后已显著优于 TF-IDF，说明\emph{语义特征}是性能增益的主要来源；而 Papex-GNN 进一步叠加异构结构建模，带来额外提升。

\begin{table}[htbp]
  \centering
  \caption{各方法在测试集上的性能（Recall@10 / MRR）}
  \label{tab:main}
  \begin{tabular}{lcc}
    \toprule
    方法 & Recall@10 & MRR \\
    \midrule
    TF-IDF            & $0.263$ & $0.141$ \\
    GCN (semantic)    & $0.352$ & $0.198$ \\
    HAN (random)      & $0.363$ & $0.205$ \\
    \textbf{Papex-GNN}& $\mathbf{0.487}$ & $\mathbf{0.271}$ \\
    \bottomrule
  \end{tabular}
\end{table}

\subsection{消融实验}
我们分别移除语义编码（替换为随机初始化）与异构边（仅保留引用边），结果见表~\ref{tab:ablation}。两项移除均导致明显退化，印证了两个模块各自的必要性。

\begin{table}[htbp]
  \centering
  \caption{消融实验（Recall@10）}
  \label{tab:ablation}
  \begin{tabular}{lc}
    \toprule
    变体 & Recall@10 \\
    \midrule
    Papex-GNN (完整)          & $0.487$ \\
    \quad 去掉语义编码        & $0.319$ \\
    \quad 去掉异构边          & $0.441$ \\
    \bottomrule
  \end{tabular}
\end{table}

\section{结论}
\section{结论}
本文提出了一种面向中文文献的图神经网络引用预测框架 Papex-GNN，通过预训练语义编码与异构图注意力结构建模的有机结合，在中文文献引用预测任务上取得了显著优于现有基线的效果。消融实验验证了两类信号各自的不可或缺性。

未来的工作将探索：(1) 引入会议、期刊等更丰富的节点类型以进一步丰富异构图结构；(2) 将框架扩展至多语言文献场景；(3) 结合大语言模型进行可解释的引用推荐。我们相信，语义与结构双轮驱动的建模思路，能够为学术知识图谱的构建与应用提供持续的动力。



% ---- 致谢 / 基金（可选，来自 papex.json 的 acknowledgments / funding） ----
% ============================================================
% 本文件由 papex-build.py 自动生成，请勿手工编辑。
% 修改请在 papex.json 中进行，随后重新构建。
% ============================================================


\papexAcknowledgments{作者感谢匿名审稿人的宝贵意见。本工作受国家自然科学基金（编号 62100000）资助的部分支持。}
\papexFunding{国家自然科学基金面上项目（62100000）；山东省重点研发计划（2026CX001）。}


% ---- 参考文献 ----
\printbibliography[heading=bibintoc,title={参考文献}]

% ---- 附录（可选，来自 papex.json 的 appendices） ----
% ============================================================
% 本文件由 papex-build.py 自动生成，请勿手工编辑。
% 修改请在 papex.json 中进行，随后重新构建。
% ============================================================


\appendix
\section{附录 A：注意力分数的闭式近似}
为降低注意力计算的开销，我们用一阶泰勒展开对式~(2) 中的注意力分数做闭式近似。记 $\mathbf{e}_{vu}^{r}=\mathbf{a}_r^{\top}[\mathbf{W}_r\mathbf{h}_v \mathbin\Vert \mathbf{W}_r\mathbf{h}_u]$，在 $\mathbf{e}_{vu}^{r}\approx 0$ 邻域内有：
\begin{equation}
  \exp(\mathrm{LeakyReLU}(\mathbf{e}_{vu}^{r})) \approx 1 + \mathbf{e}_{vu}^{r}.
\end{equation}
代入归一化项可得近似注意力权重 $\tilde{\alpha}_{vu}^{r}$，其计算复杂度由指数运算降为线性，在大规模图上可带来约 $1.8\times$ 的加速，且精度损失小于 $0.5\%$。



\end{document}
